\chapter{Introduction}

This blueprint accompanies a \emph{sorry-free} Lean~4 + Mathlib formalization of the
\textbf{Oseledets multiplicative ergodic theorem} (MET) and a broad layer of companion
results. It documents the mathematical content of the development and the dependency
structure of its proof: every node carries a \texttt{\textbackslash lean} annotation pointing at the
corresponding Lean declaration, and a green checkmark records that the statement (and,
where shown, its proof) is fully formalized. The API documentation generated by
\texttt{doc-gen4} is linked from each node.

\section{Setting}

Throughout, \((X, \mu)\) is a probability space and \(T : X \to X\) is an ergodic
measure-preserving transformation. A \emph{linear cocycle} of dimension \(d\) is generated
by a measurable map
\[
  A : X \to \mathrm{Mat}_{d}(\R), \qquad \det(A\,x) \neq 0 \text{ for all } x,
\]
whose \(n\)-step product along the orbit of \(x\) is
\[
  A^{(n)}(x) \;=\; A(T^{n-1}x)\cdots A(Tx)\,A(x), \qquad A^{(0)}(x) = I .
\]
We impose the one-sided integrability condition
\[
  \log^{+}\norm{A}, \ \log^{+}\norm{A^{-1}} \in L^{1}(\mu),
\]
where matrices act on \(\mathrm{EuclideanSpace}\ \R\ (\mathrm{Fin}\ d)\) so that
\(\norm{\cdot}\) is the \(L^2\) operator norm (submultiplicative). For \(v \neq 0\) the
\emph{Lyapunov exponent} in the direction \(v\) is the growth rate
\[
  \bar\lambda(x, v) \;=\; \limsup_{n\to\infty} \tfrac1n \log \norm{A^{(n)}(x)\, v}.
\]

\section{The three headline theorems}

The development proves three principal theorems, each formalized sorry-free and verified
(by a guarded axiom audit) to depend only on the standard axioms
\(\{\texttt{propext}, \texttt{Classical.choice}, \texttt{Quot.sound}\}\).

\begin{itemize}
  \item \textbf{One-sided MET (filtration form),} \lean{Oseledets.oseledets_filtration}: there
    are finitely many distinct Lyapunov exponents \(\lambda_1 > \cdots > \lambda_k\) and,
    for \(\mu\)-a.e.\ \(x\), a strictly decreasing, \(A\)-equivariant, measurable filtration
    \[
      \mathrm{EuclideanSpace}\ \R\ (\mathrm{Fin}\ d) = V^0_x \supsetneq V^1_x \supsetneq
      \cdots \supsetneq V^k_x = \{0\}
    \]
    along which \(\tfrac1n \log\norm{A^{(n)}(x)\,v} \to \lambda_i\) for every
    \(v \in V^i_x \setminus V^{i+1}_x\). This is the central result; it is proved in
    Chapter~\ref{chap:met}.
  \item \textbf{Two-sided splitting,} \lean{Oseledets.oseledets_splitting}: when both the
    forward and backward filtrations are available, they are transverse and the space
    splits a.e.\ into an \(A\)-equivariant direct sum of \emph{Oseledets subspaces}
    \(\R^d = \bigoplus_i E^i_x\) on which the cocycle grows at the exact two-sided rate
    \(\lambda_i\). Chapter~\ref{chap:twosided}.
  \item \textbf{Continuous-flow MET,} \lean{Oseledets.oseledets_flow}: the analogue for a
    continuous-time linear cocycle over a measurable flow, obtained by reduction to the
    time-one map together with a between-times sandwich estimate. Chapter~\ref{chap:continuous}.
\end{itemize}

\section{Structure of the proof}

The proof proceeds in layers, mirrored by the chapters of this blueprint.
Chapter~\ref{chap:cocycle} sets up the cocycle, the operator norm and its measurability,
and the Furstenberg--Kesten extremal exponents. Chapter~\ref{chap:ergodic} develops the
ergodic-theoretic engine: the maximal ergodic inequality, the pointwise Birkhoff theorem,
and---crucially---Kingman's subadditive ergodic theorem, which converts the subadditivity
of the log-norms into almost-everywhere limits. Chapter~\ref{chap:lyapunov} introduces the
Lyapunov exponent as a measurable function, the Lyapunov spectrum, and the limsup
filtration together with the measurability of its subspaces. Chapter~\ref{chap:met}
assembles the one-sided theorem through the Oseledets limit, the spectral upper bound and
determinant squeeze, the spectral identification of the limsup filtration, and a top-gap
envelope induction. The remaining chapters record the companion results
(Chapter~\ref{chap:corollaries}), the two-sided splitting (Chapter~\ref{chap:twosided}),
and the continuous-flow theorem (Chapter~\ref{chap:continuous}).
